\documentclass{amsart}
% For dealing with references we use the comment environment
\usepackage{verbatim}
\newenvironment{reference}{\comment}{\endcomment}
%\newenvironment{reference}{}{}
\newenvironment{slogan}{\comment}{\endcomment}
\newenvironment{history}{\comment}{\endcomment}

% For commutative diagrams we use Xy-pic
\usepackage[all]{xy}

% We use 2cell for 2-commutative diagrams.
\xyoption{2cell}
\UseAllTwocells

% We use multicol for the list of chapters between chapters
\usepackage{multicol}

% This is generally recommended for better output
\usepackage{lmodern}
\usepackage[T1]{fontenc}

% For cross-file-references

% Package for hypertext links:
\usepackage{hyperref}

% For any local file, say "hello.tex" you want to link to please

% Theorem environments.
%
\theoremstyle{plain}
\newtheorem{theorem}[subsection]{Theorem}
\newtheorem{proposition}[subsection]{Proposition}
\newtheorem{lemma}[subsection]{Lemma}

\theoremstyle{definition}
\newtheorem{definition}[subsection]{Definition}
\newtheorem{example}[subsection]{Example}
\newtheorem{exercise}[subsection]{Exercise}
\newtheorem{situation}[subsection]{Situation}

\theoremstyle{remark}
\newtheorem{remark}[subsection]{Remark}
\newtheorem{remarks}[subsection]{Remarks}

\numberwithin{equation}{subsection}

% Macros
%
\def\lim{\mathop{\mathrm{lim}}\nolimits}
\def\colim{\mathop{\mathrm{colim}}\nolimits}
\def\Spec{\mathop{\mathrm{Spec}}}
\def\Hom{\mathop{\mathrm{Hom}}\nolimits}
\def\Ext{\mathop{\mathrm{Ext}}\nolimits}
\def\SheafHom{\mathop{\mathcal{H}\!\mathit{om}}\nolimits}
\def\SheafExt{\mathop{\mathcal{E}\!\mathit{xt}}\nolimits}
\def\Sch{\mathit{Sch}}
\def\Mor{\mathop{\mathrm{Mor}}\nolimits}
\def\Ob{\mathop{\mathrm{Ob}}\nolimits}
\def\Sh{\mathop{\mathit{Sh}}\nolimits}
\def\NL{\mathop{N\!L}\nolimits}
\def\CH{\mathop{\mathrm{CH}}\nolimits}
\def\proetale{{pro\text{-}\acute{e}tale}}
\def\etale{{\acute{e}tale}}
\def\QCoh{\mathit{QCoh}}
\def\Ker{\mathop{\mathrm{Ker}}}
\def\Im{\mathop{\mathrm{Im}}}
\def\Coker{\mathop{\mathrm{Coker}}}
\def\Coim{\mathop{\mathrm{Coim}}}

% Boxtimes
%
\DeclareMathSymbol{\boxtimes}{\mathbin}{AMSa}{"02}

%
% Macros for moduli stacks/spaces
%
\def\QCohstack{\mathcal{QC}\!\mathit{oh}}
\def\Cohstack{\mathcal{C}\!\mathit{oh}}
\def\Spacesstack{\mathcal{S}\!\mathit{paces}}
\def\Quotfunctor{\mathrm{Quot}}
\def\Hilbfunctor{\mathrm{Hilb}}
\def\Curvesstack{\mathcal{C}\!\mathit{urves}}
\def\Polarizedstack{\mathcal{P}\!\mathit{olarized}}
\def\Complexesstack{\mathcal{C}\!\mathit{omplexes}}
% \Pic is the operator that assigns to X its picard group, usage \Pic(X)
% \Picardstack_{X/B} denotes the Picard stack of X over B
% \Picardfunctor_{X/B} denotes the Picard functor of X over B
\def\Pic{\mathop{\mathrm{Pic}}\nolimits}
\def\Picardstack{\mathcal{P}\!\mathit{ic}}
\def\Picardfunctor{\mathrm{Pic}}
\def\Deformationcategory{\mathcal{D}\!\mathit{ef}}

\setlength{\emergencystretch}{3em}
\begin{document}
\title[Stacks Project, Tag 06KT]{402 --- Rim gluing for quotients of smooth groupoids in functors}
\maketitle
\noindent Copyright (C) 2005--2025 Johan de Jong. Stacks Project authors. Source: \href{https://stacks.math.columbia.edu/tag/06KT}{Tag 06KT}.

\noindent Editorial difficulty note (not author text). Difficulty H2: The proof constructs an explicit quotient fibre square by lifting an arrow through the smooth source map. It then checks both commutativity and the universal property, reconstructing the unique mediating arrow from the Rim gluing property of the arrow functor. This is effective two-categorical gluing for a quotient, rather than formal preservation of a property..

\noindent Editorial provenance note (not author text). History: excerpt from revision \texttt{a0444\allowbreak 6e57e\allowbreak c1fbc\allowbreak 252a8\allowbreak 71afc\allowbreak ec775\allowbreak 2fb28\allowbreak 07b14}, formal-defos.tex, lines 6496--6572. Reference presentation was adapted; original-author context and core support blocks were retained or added. The mathematical wording is unchanged. Licensed under GNU FDL 1.2 or later; no invariant sections or cover texts. See ../licenses/COPYING. Context blocks reproduce author definitions and supporting results; they are not additional counted problems.

\begin{definition}
\label{definition-CLambda}
Let $\Lambda$ be a Noetherian ring and let $\Lambda \to k$ be a finite
ring map where $k$ is a field. We define {\it $\mathcal{C}_\Lambda$} to be
the category with
\begin{enumerate}
\item objects are pairs $(A, \varphi)$ where $A$ is an Artinian local
$\Lambda$-algebra and where $\varphi : A/\mathfrak m_A \to k$ is a
$\Lambda$-algebra isomorphism, and
\item morphisms $f : (B, \psi) \to (A, \varphi)$ are local $\Lambda$-algebra
homomorphisms such that $\varphi \circ (f \bmod \mathfrak m) = \psi$.
\end{enumerate}
We say we are in the {\it classical case} if $\Lambda$ is a Noetherian
complete local ring and $k$ is its residue field.
\end{definition}

\section{Notation and Conventions}
\label{section-notations-conventions}

\noindent
A ring is commutative with $1$. The  maximal ideal of a local ring $A$
is denoted by $\mathfrak{m}_A$. The set of positive integers is denoted
by $\mathbf{N} = \{1, 2, 3, \ldots\}$. If $U$ is an object of a
category $\mathcal{C}$, we denote by $\underline{U}$
the functor
$\Mor_\mathcal{C}(U, -): \mathcal{C} \to \textit{Sets}$, see
Remarks \ref{remarks-cofibered-groupoids} (\ref{item-definition-yoneda}).
Warning: this may conflict with the notation in other chapters where we
sometimes use $\underline{U}$ to denote $h_U(-) = \Mor_\mathcal{C}(-, U)$.

\medskip\noindent
Throughout this chapter $\Lambda$ is a Noetherian ring and
$\Lambda \to k$ is a finite ring map from $\Lambda$ to a field.
The kernel of this map is denoted $\mathfrak m_\Lambda$ and the
image $k' \subset k$. It turns out that $\mathfrak m_\Lambda$ is
a maximal ideal, $k' = \Lambda/\mathfrak m_\Lambda$ is a field, and
the extension $k/k'$ is finite. See discussion surrounding
(\href{https://stacks.math.columbia.edu/tag/06S2}{equation k prime (Tag 06S2)}).




\begin{definition}
\label{definition-completion-CLambda}
Let $\Lambda$ be a Noetherian ring and let $\Lambda \to k$ be a finite
ring map where $k$ is a field. We define {\it $\widehat{\mathcal{C}}_\Lambda$}
to be the category with
\begin{enumerate}
\item objects are pairs $(R, \varphi)$ where $R$ is a Noetherian complete
local $\Lambda$-algebra and where $\varphi : R/\mathfrak m_R \to k$ is a
$\Lambda$-algebra isomorphism, and
\item morphisms $f : (S, \psi) \to (R, \varphi)$ are local $\Lambda$-algebra
homomorphisms such that $\varphi \circ (f \bmod \mathfrak m) = \psi$.
\end{enumerate}
\end{definition}

\begin{definition}
\label{definition-category-cofibred-groupoids}
Let $\mathcal{C}$ be a category.  A {\it category cofibered in groupoids over
$\mathcal{C}$} is a category $\mathcal{F}$ equipped with a functor
$p: \mathcal{F} \to \mathcal{C}$ such that $\mathcal{F}^{opp}$ is a category
fibered in groupoids over $\mathcal{C}^{opp}$ via
$p^{opp}: \mathcal{F}^{opp} \to \mathcal{C}^{opp}$.
\end{definition}

\begin{definition}
\label{definition-predeformation-category}
A {\it predeformation category} $\mathcal{F}$ is a category cofibered
in groupoids over $\mathcal{C}_\Lambda$ such that $\mathcal{F}(k)$ is
equivalent to a category with a single object and a single morphism,
i.e., $\mathcal{F}(k)$ contains at least one object and there is a
unique morphism between any two objects. A {\it morphism of predeformation
categories} is a morphism of categories cofibered in groupoids over
$\mathcal{C}_\Lambda$.
\end{definition}

\begin{definition}
\label{definition-formal-objects}
Let $\mathcal{F}$ be a category cofibered in groupoids over
$\mathcal{C}_\Lambda$. The {\it category $\widehat{\mathcal{F}}$ of formal
objects of  $\mathcal{F}$} is the category with the following objects and
morphisms.
\begin{enumerate}
\item A {\it formal object $\xi = (R, \xi_n, f_n)$ of $\mathcal{F}$}
consists of an object $R$ of $\widehat{\mathcal{C}}_\Lambda$, and a collection
indexed by $n \in \mathbf{N}$ of objects $\xi_n$ of
$\mathcal{F}(R/\mathfrak m_R^n)$ and morphisms
$f_n : \xi_{n + 1} \to \xi_n$ lying over the projection
$R/\mathfrak m_R^{n + 1} \to R/\mathfrak m_R^n$.
\item Let $\xi = (R, \xi_n, f_n)$ and $\eta = (S, \eta_n, g_n)$ be
formal objects of $\mathcal{F}$.  A {\it morphism $a : \xi \to \eta$ of
formal objects} consists of a map $a_0 : R \to S$ in
$\widehat{\mathcal{C}}_\Lambda$ and a collection $a_n : \xi_n \to \eta_n$
of morphisms of $\mathcal{F}$ lying over
$R/\mathfrak m_R^n \to S/\mathfrak m_S^n$,
such that for every $n$ the diagram
$$
\xymatrix{
\xi_{n + 1} \ar[r]_{f_n} \ar[d]_{a_{n + 1}} & \xi_n \ar[d]^{a_n} \\
\eta_{n + 1} \ar[r]^{g_n} & \eta_n
}
$$
commutes.
\end{enumerate}
\end{definition}

\begin{definition}
\label{definition-versal}
Let $\mathcal{F}$ be a category cofibered in groupoids.  Let $\xi$ be a formal
object of $\mathcal{F}$ lying over $R \in \Ob(\widehat{\mathcal{C}}_\Lambda)$.
We say $\xi$ is {\it versal} if the corresponding morphism
$\underline{\xi}: \underline{R}|_{\mathcal{C}_\Lambda} \to \mathcal{F}$
of Remark \href{https://stacks.math.columbia.edu/tag/06HC}{remark formal objects yoneda (Tag 06HC)} is smooth.
\end{definition}

\begin{definition}
\label{definition-S1-S2}
Let $\mathcal{F}$ be a category cofibered in groupoids over $\mathcal
C_\Lambda$. We define {\it conditions (S1) and (S2)}
on $\mathcal{F}$ as follows:
\begin{enumerate}
\item[(S1)] Every diagram in $\mathcal{F}$
$$
\vcenter{
\xymatrix{
           & x_2 \ar[d] \\
x_1 \ar[r] & x
}
}
\quad\text{lying over}\quad
\vcenter{
\xymatrix{
           & A_2 \ar[d] \\
A_1 \ar[r] & A
}
}
$$
in $\mathcal{C}_\Lambda$ with $A_2 \to A$ surjective can be completed
to a commutative diagram
$$
\vcenter{
\xymatrix{
y \ar[r] \ar[d] & x_2 \ar[d] \\
x_1 \ar[r]      & x
}
}
\quad\text{lying over}\quad
\vcenter{
\xymatrix{
A_1 \times_A A_2 \ar[r] \ar[d] & A_2 \ar[d] \\
A_1 \ar[r]      & A.
}
}
$$
\item[(S2)]
The condition of (S1) holds for diagrams in $\mathcal{F}$ lying over
a diagram in $\mathcal{C}_\Lambda$ of the form
$$
\xymatrix{
          & k[\epsilon] \ar[d] \\
A  \ar[r] & k.
}
$$
Moreover, if we have two commutative diagrams in $\mathcal{F}$
$$
\vcenter{
\xymatrix{
y \ar[r]_c \ar[d]_a & x_\epsilon \ar[d]^e \\
x \ar[r]^d          & x_0
}
}
\quad\text{and}\quad
\vcenter{
\xymatrix{
y' \ar[r]_{c'} \ar[d]_{a'} & x_\epsilon \ar[d]^e \\
x \ar[r]^d                 & x_0
}
}
\quad\text{lying over}\quad
\vcenter{
\xymatrix{
A \times_k k[\epsilon] \ar[r] \ar[d] & k[\epsilon] \ar[d] \\
A  \ar[r] & k
}
}
$$
then there exists a morphism $b : y \to y'$ in
$\mathcal{F}(A \times_k k[\epsilon])$ such that $a = a' \circ b$.
\end{enumerate}
\end{definition}

\begin{definition}
\label{definition-RS}
Let $\mathcal{F}$ be a category cofibered in groupoids over $\mathcal
C_\Lambda$.  We say that $\mathcal{F}$ satisfies {\it condition (RS)}
if for every diagram in $\mathcal{F}$
$$
\vcenter{
\xymatrix{
           & x_2 \ar[d] \\
x_1 \ar[r] & x
}
}
\quad\text{lying over}\quad
\vcenter{
\xymatrix{
           & A_2 \ar[d] \\
A_1 \ar[r] & A
}
}
$$
in $\mathcal{C}_\Lambda$ with $A_2 \to A$ surjective, there exists a
fiber product $x_1 \times_x x_2$ in $\mathcal{F}$ such that the diagram
$$
\vcenter{
\xymatrix{
x_1 \times_x x_2 \ar[r] \ar[d] & x_2 \ar[d] \\
x_1 \ar[r]      & x
}
}
\quad\text{lies over}\quad
\vcenter{
\xymatrix{
A_1 \times_A A_2 \ar[r] \ar[d] & A_2 \ar[d] \\
A_1 \ar[r]      & A.
}
}
$$
\end{definition}

\begin{definition}
\label{definition-groupoid-in-functors}
Let $\mathcal{C}$ be a category. The
{\it category of groupoids in functors on $\mathcal{C}$}
is the category with the following objects and morphisms.
\begin{enumerate}
\item Objects: A {\it groupoid in functors on $\mathcal{C}$} is a quintuple
$(U, R, s, t, c)$ where $U, R : \mathcal{C} \to \textit{Sets}$ are
functors and $s, t : R \to U$ and $c : R \times_{s, U, t} R \to R$
are morphisms with the following property: For any object $T$ of $\mathcal{C}$,
the quintuple
$$
(U(T), R(T), s, t, c)
$$
is a groupoid category.
\item Morphisms: A {\it morphism $(U, R, s, t, c) \to (U', R', s', t', c')$ of
groupoids in functors on $\mathcal{C}$} consists of morphisms $U \to U'$
and $R \to R'$ with the following property: For any object $T$ of
$\mathcal{C}$, the induced maps $U(T) \to U'(T)$ and
$R(T) \to R'(T)$ define a functor between groupoid categories
$$
(U(T), R(T), s, t, c) \to (U'(T), R'(T), s', t', c').
$$
\end{enumerate}
\end{definition}

\begin{definition}
\label{definition-quotient}
Let $(U, R, s, t, c)$ be a groupoid in functors on a category $\mathcal{C}$.
\begin{enumerate}
\item The assignment $T \mapsto  (U(T), R(T), s, t, c)$ determines a functor
$\mathcal{C} \to \textit{Groupoids}$. The {\it quotient category
cofibered in groupoids $[U/R] \to \mathcal{C}$} is the category
cofibered in groupoids over $\mathcal{C}$ associated to this functor (as in
Remarks \ref{remarks-cofibered-groupoids}
(\ref{item-construction-associated-cofibered-groupoid})).
\item The {\it quotient morphism $U \to [U/R]$} is the morphism of
categories cofibered in groupoids over $\mathcal{C}$ defined by the
rules
\begin{enumerate}
\item $x \in U(T)$ maps to the object $(T, x) \in \Ob([U/R](T))$, and
\item $x \in U(T)$ and $f : T \to T'$ give rise to the morphism
$(f, \text{id}_{U(f)(x)}): (T, x) \to (T, U(f)(x))$ lying over
$f : T \to T'$.
\end{enumerate}
\end{enumerate}
\end{definition}

\begin{definition}
\label{definition-smooth-groupoid-in-functors}
Let $(U, R, s, t, c)$ be a groupoid in functors on $\mathcal{C}_\Lambda$. We
say $(U, R, s, t, c)$ is {\it smooth} if $s, t: R \to U$ are smooth.
\end{definition}

\begin{definition}
\label{definition-smooth-morphism}
Let $\varphi : \mathcal{F} \to \mathcal{G}$ be a morphism of categories
cofibered in groupoids over $\mathcal{C}_\Lambda$.  We say  $\varphi$ is
{\it smooth} if it satisfies the following condition: Let $B \to A$ be
a surjective ring map in $\mathcal{C}_\Lambda$.  Let $y \in
\Ob(\mathcal{G}(B)), x \in \Ob(\mathcal{F}(A))$, and $y
\to \varphi(x)$ be a morphism lying over $B \to A$.  Then there
exists $x' \in \Ob(\mathcal{F}(B))$, a morphism $x' \to x$
lying over $B \to A$, and a morphism $\varphi(x') \to y$ lying
over $\text{id}: B \to B$, such that the diagram
$$
\xymatrix{
\varphi(x') \ar[r] \ar[dr] & y \ar[d] \\
& \varphi(x)
}
$$
commutes.
\end{definition}

\begin{remarks}
\label{remarks-cofibered-groupoids}
Everything about categories fibered in groupoids translates directly to the
cofibered setting. The following remarks are meant to fix notation.
Let $\mathcal{C}$ be a category.
\begin{enumerate}
\item We often omit the functor $p: \mathcal{F} \to \mathcal{C}$ from the
notation.
\item The fiber category over an object $U$ in $\mathcal{C}$ is denoted by
$\mathcal{F}(U)$. Its objects are those of $\mathcal{F}$ lying over $U$ and its
morphisms are those of $\mathcal{F}$ lying over $\text{id}_U$.
If $x, y$  are objects of $\mathcal{F}(U)$, we sometimes write
$\Mor_U(x, y)$ for $\Mor_{\mathcal{F}(U)}(x, y)$.
\item The fibre categories $\mathcal{F}(U)$ are groupoids, see
Categories, Lemma \href{https://stacks.math.columbia.edu/tag/003V}{lemma fibred groupoids (Tag 003V)}.
Hence the morphisms in $\mathcal{F}(U)$ are all isomorphisms.
We sometimes write $\text{Aut}_U(x)$ for $\Mor_{\mathcal{F}(U)}(x, x)$.
\item
\label{item-pushforward}
Let $\mathcal{F}$ be a category cofibered in groupoids over
$\mathcal{C}$, let $f: U \to V$ be a morphism in $\mathcal{C}$, and
let $x \in \Ob(\mathcal{F}(U))$.
A {\it pushforward} of $x$ along $f$ is a morphism
$x \to y$ of $\mathcal{F}$ lying over $f$. A pushforward
is unique up to unique isomorphism (see the discussion following
Categories, Definition \ref{categories-definition-cartesian-over-C}).
We sometimes write $x \to f_*x$ for ``the'' pushforward of $x$
along $f$.
\item A {\it choice of pushforwards for $\mathcal{F}$} is the choice of
a pushforward of $x$ along $f$ for every pair $(x, f)$ as above. We can make
such a choice of pushforwards for $\mathcal{F}$ by the axiom of choice.
\item Let $\mathcal{F}$ be a category cofibered in groupoids over
$\mathcal{C}$. Given a choice of pushforwards for $\mathcal{F}$, there
is an associated pseudo-functor $\mathcal{C} \to \textit{Groupoids}$.
We will never use this construction so we give no details.
\item
\label{item-cofibered-morphism}
A morphism of categories cofibered in groupoids over $\mathcal{C}$ is a
functor commuting with the projections to $\mathcal{C}$. If $\mathcal{F}$
and $\mathcal{F}'$ are categories cofibered in groupoids over
$\mathcal{C}$, we denote the morphisms from $\mathcal{F}$ to $\mathcal{F}'$
by $\Mor_\mathcal{C}(\mathcal{F}, \mathcal{F}')$.
\item
\label{item-definition-cofibered-groupoids-2-category}
Categories cofibered in groupoids form a $(2, 1)$-category
$\text{Cof}(\mathcal{C})$. Its 1-morphisms are the morphisms described in
(\ref{item-cofibered-morphism}). If $p : \mathcal{F} \to C$ and
$p': \mathcal{F}' \to \mathcal{C}$ are categories cofibered in groupoids
and $\varphi, \psi : \mathcal{F} \to \mathcal{F}'$ are $1$-morphisms, then
a 2-morphism $t : \varphi \to \psi$ is a morphism of functors such that
$p'(t_x) = \text{id}_{p(x)}$ for all $x \in \Ob(\mathcal{F})$.
\item
\label{item-construction-associated-cofibered-groupoid}
Let $F : \mathcal{C} \to \textit{Groupoids}$ be a functor. There
is a category cofibered in groupoids $\mathcal{F} \to \mathcal{C}$
associated to $F$ as follows. An object of $\mathcal{F}$ is a pair $(U, x)$
where $U \in \Ob(\mathcal{C})$ and $x \in \Ob(F(U))$. A
morphism $(U, x) \to (V, y)$ is a pair $(f, a)$ where
$f \in \Mor_\mathcal{C}(U, V)$ and
$a \in \Mor_{F(V)}(F(f)(x), y)$.
The functor $\mathcal{F} \to \mathcal{C}$ sends $(U, x)$ to $U$. See
Categories, Section \href{https://stacks.math.columbia.edu/tag/0048}{section presheaves groupoids (Tag 0048)}.
\item
\label{item-associated-functor-isomorphism-classes}
Let $\mathcal{F}$ be cofibered in groupoids over $\mathcal{C}$.
For $U \in \Ob(\mathcal{C})$ set $\overline{\mathcal{F}}(U)$ equal to
the set of isomorphisms classes of the category $\mathcal{F}(U)$.
If $f : U \to V$ is a morphism of $\mathcal{C}$, then we obtain a
map of sets $\overline{\mathcal{F}}(U) \to \overline{\mathcal{F}}(V)$ by
mapping the isomorphism class of $x$ to the isomorphism class of a pushforward
$f_*x$ of $x$ see (\ref{item-pushforward}). Then
$\overline{\mathcal{F}} : \mathcal{C} \to \textit{Sets}$ is a
functor. Similarly, if $\varphi : \mathcal{F} \to \mathcal{G}$ is a
morphism of cofibered categories, we denote by
$\overline{\varphi}: \overline{\mathcal{F}} \to  \overline{\mathcal{G}}$
the associated morphism of functors.
\item
\label{item-convention-cofibered-sets}
Let $F: \mathcal{C} \to \textit{Sets}$ be a functor. We can think of a
set as a discrete category, i.e., as a groupoid with only identity morphisms.
Then the construction (\ref{item-construction-associated-cofibered-groupoid})
associates to $F$ a category cofibered in sets. This defines a fully
faithful embedding of the category of functors $\mathcal{C} \to \textit{Sets}$
to the category of categories cofibered in groupoids over $\mathcal{C}$.
We identify the category of functors with its image under this embedding.
Hence if $F : \mathcal{C} \to \textit{Sets}$ is a functor, we denote the
associated category cofibered in sets also by $F$; and if
$\varphi : F \to G$ is a morphism of functors, we denote still by $\varphi$
the corresponding morphism of categories cofibered in sets, and vice-versa.
See Categories, Section \href{https://stacks.math.columbia.edu/tag/0042}{section fibred in sets (Tag 0042)}.
\item
\label{item-definition-yoneda}
Let $U$ be an object of $\mathcal{C}$.  We write $\underline{U}$ for the
functor
$\Mor_\mathcal{C}(U, -): \mathcal{C} \to
\textit{Sets}$.  This defines a fully faithful embedding of $\mathcal
C^{opp}$ into the category of functors $\mathcal{C} \to
\textit{Sets}$. Hence, if $f : U \to V$ is a morphism, we are
justified in denoting still by $f$ the induced morphism $\underline{V}
\to \underline{U}$, and vice-versa.
\item
\label{item-fibre-product}
Fiber products of categories cofibered in groupoids: If $\mathcal{F}
\to \mathcal{H}$ and $\mathcal{G} \to \mathcal{H}$ are morphisms
of categories cofibered in groupoids over $\mathcal{C}_\Lambda$, then a
construction of their 2-fiber product is given by the construction for their
2-fiber product as categories over $\mathcal{C}_\Lambda$, as described in
Categories, Lemma \href{https://stacks.math.columbia.edu/tag/0040}{lemma 2 product categories over C (Tag 0040)}.
\item
\label{item-product}
Products of categories cofibered in groupoids: If $\mathcal{F}$ and
$\mathcal{G}$ are categories cofibered in groupoids over
$\mathcal{C}_\Lambda$ then their product is defined to be the $2$-fiber product
$\mathcal{F} \times_{\mathcal{C}_\Lambda} \mathcal{G}$ as described in
Categories, Lemma \href{https://stacks.math.columbia.edu/tag/0040}{lemma 2 product categories over C (Tag 0040)}.
\item
\label{item-definition-restricting-base-category}
Restricting the base category: Let $p : \mathcal{F} \to \mathcal{C}$ be a
category cofibered in groupoids, and let $\mathcal{C}'$ be a full
subcategory of $\mathcal{C}$. The restriction $\mathcal{F}|_{\mathcal{C}'}$
is the full subcategory of $\mathcal{F}$ whose objects lie over
objects of $\mathcal{C}'$. It is a category cofibered in groupoids via
the functor
$p|_{\mathcal{C}'}: \mathcal{F}|_{\mathcal{C}'} \to \mathcal{C}'$.
\end{enumerate}
\end{remarks}

\begin{definition}
\label{categories-definition-cartesian-over-C}
Let $\mathcal{C}$ be a category.
Let $p : \mathcal{S} \to \mathcal{C}$ be a category over $\mathcal{C}$.
A {\it strongly cartesian morphism}, or more precisely a
{\it strongly $\mathcal{C}$-cartesian morphism} is a
morphism $\varphi : y \to x$ of $\mathcal{S}$ such that
for every $z \in \Ob(\mathcal{S})$ the map
$$
\Mor_\mathcal{S}(z, y)
\longrightarrow
\Mor_\mathcal{S}(z, x)
\times_{\Mor_\mathcal{C}(p(z), p(x))}
\Mor_\mathcal{C}(p(z), p(y)),
$$
given by $\psi \longmapsto (\varphi \circ \psi, p(\psi))$
is bijective.
\end{definition}

\noindent
Explicitly, $p: \mathcal{F} \to \mathcal{C}$ is cofibered in groupoids if
the following two conditions hold:
\begin{enumerate}
\item For every morphism $f: U \to V$ in $\mathcal{C}$ and every object
$x$ lying over $U$, there is a morphism $x \to y$ of $\mathcal{F}$ lying
over $f$.
\item For every pair of morphisms $a: x \to y$ and $b: x \to z$
of $\mathcal{F}$ and any morphism $f: p(y) \to p(z)$ such that $p(b) = f
\circ p(a)$, there exists a unique morphism $c: y \to z$ of $\mathcal
F$ lying over $f$ such that $b = c \circ a$.
\end{enumerate}


\medskip\noindent
We denote by $k[\epsilon]$ the ring of dual numbers over $k$.  More
generally, for a $k$-vector space $V$, we denote by $k[V]$ the $k$-algebra
whose underlying vector space is $k \oplus V$ and whose multiplication is given
by $(a, v) \cdot (a', v') = (aa', av' + a'v)$.  When $V = k$, $k[V]$ is the ring
of dual numbers over $k$.  For any finite dimensional $k$-vector space $V$
the ring $k[V]$ is in $\mathcal{C}_\Lambda$.


\begin{lemma}
\label{lemma-smooth-RS-groupoid-in-functors-quotient}
Let $(U, R, s, t, c)$ be a smooth groupoid in functors on $\mathcal{C}_\Lambda$.
Assume $U$ and $R$ satisfy (RS). Then $[U/R]$ satisfies (RS).
\end{lemma}

\begin{proof}
Let
$$
\xymatrix{
                           &     (A_2, x_2) \ar[d]^{(f_2, a_2)} \\
(A_1, x_1) \ar[r]^{(f_1, a_1)} &     (A, x)
}
$$
be a diagram in $[U/R]$ such that $f_2: A_2 \to A$ is surjective. The
notation is as in
Remarks \ref{remarks-cofibered-groupoids}
(\ref{item-construction-associated-cofibered-groupoid}).
Hence $f_1: A_1 \to A, f_2: A_2 \to A$ are maps in
$\mathcal{C}_\Lambda$, $x \in U(A)$, $x_1 \in U(A_1)$, $x_2 \in U(A_2)$,
and $a_1, a_2 \in R(A)$ with $s(a_1) = U(f_1)(x_1)$,
$t(a_1) = x$ and $s(a_2) = U(f_2)(x_2)$, $t(a_2) = x$.
We construct a fiber product lying over $A_1 \times_A A_2$
for this diagram in $[U/R]$ as follows.

\medskip\noindent
Let $a = c(i(a_1), a_2)$, where $i: R \to R$ is the inverse morphism.
Then $a \in R(A)$, $x_2 \in U(A_2)$ and $s(a) = U(f_2)(x_2)$.
Hence an element $(a, x_2) \in R(A) \times_{s, U(A), U(f_2)} U(A_2)$.
By smoothness of $s : R \to U$ there is an element
$\widetilde{a} \in R(A_2)$ with $R(f_2)(\widetilde{a}) = a$ and
$s(\widetilde{a}) = x_2$. In particular
$U(f_2)(t(\widetilde{a})) = t(a) = U(f_1)(x_1)$. Thus $x_1$ and
$t(\widetilde{a})$ define an element
$$
(x_1, t(\widetilde{a})) \in U(A_1) \times_{U(A)} U(A_2).
$$
By the assumption that $U$ satisfies (RS), we have an identification
$U(A_1) \times_{U(A)} U(A_2) = U(A_1 \times_A A_2)$. Let us denote
$x_1 \times t(\widetilde{a}) \in U(A_1 \times_A A_2)$ the element
corresponding to $(x_1, t(\widetilde{a})) \in U(A_1) \times_{U(A)} U(A_2)$.
Let $p_1, p_2$ be the projections of $A_1 \times_A A_2$.  We claim
$$
\xymatrix{
(A_1 \times_A A_2, x_1 \times t(\widetilde{a})) \ar[d]_{(p_1, e(x_1))}
\ar[rr]_-{(p_2, i(\widetilde{a}))} & & (A_2, x_2) \ar[d]^{(f_2, a_2)} \\
(A_1, x_1) \ar[rr]^{(f_1, a_1)} & & (A, x)
}
$$
is a fiber square in $[U/R]$. (Note $e: U \to R$ denotes the identity.)

\medskip\noindent
The diagram is commutative because
$c(a_2, R(f_2)(i(\widetilde{a}))) = c(a_2, i(a)) = a_1$.
To check it is a fiber square, let
$$
\xymatrix{
(B, z) \ar[d]_{(g_1, b_1)} \ar[rr]_{(g_2, b_2)} & & (A_2, x_2)
\ar[d]^{(f_2, a_2)} \\
(A_1, x_1) \ar[rr]^{(f_1, a_1)} & & (A, x)
}
$$
be a commutative diagram in $[U/R]$. We will show there is a unique morphism
$(g, b) : (B, z) \to (A_1 \times_A A_2, x_1 \times t(\widetilde{a}))$
compatible with the morphisms to $(A_1, x_1)$ and $(A_2, x_2)$.
We must take $g = (g_1, g_2) : B \to A_1 \times_A A_2$.
Since by assumption $R$ satisfies (RS), we have an identification
$R(A_1 \times_A A_2) = R(A_1) \times_{R(A)} R(A_2)$.
Hence we can write $b = (b'_1, b'_2)$ for some
$b'_1 \in R(A_1)$, $b'_2 \in R(A_2)$ which agree in $R(A)$.
Then
$((g_1, g_2), (b'_1, b'_2)) : (B, z) \to
(A_1 \times_A A_2, x_1 \times t(\widetilde{a}))$
will commute with the projections if and only if
$b'_1 = b_1$ and $b'_2 = c(\widetilde{a}, b_2)$ proving unicity and
existence.
\end{proof}
\end{document}
